% Source - https://tex.stackexchange.com/a/330103
% Posted by user2574, modified by community. See post 'Timeline' for change history
% Retrieved 2026-02-18, License - CC BY-SA 4.0
\documentclass[a4paper, 12pt]{article}
\usepackage[margin=1in]{geometry} % Adjusts standard page margins
\usepackage{caption}
\usepackage{subcaption}
\usepackage{graphicx}

\usepackage[T1]{fontenc}
\usepackage[scaled]{helvet} % Loads Helvetica
\scriptsize 


\renewcommand{\familydefault}{\sfdefault} % Sets Helvetica as default




\begin{document}
\noindent\begin{center}
\begin{figure}
%\hspace{-1.2em}
 %\makebox[\textwidth][l]{\scriptsize{Numbers of DTU features}}
 
 
  \begin{subfigure}[t]{0.2\textwidth}
        \centering
        \llap{\raisebox{0mm}{\textbf{\tiny a }}}%
               \vspace{-2mm}
      \tiny  Transcript proportions (ACTG1)
    %    \makebox[\linewidth][l]{\tiny \textbf{a} Proportions}
        \includegraphics[width=\linewidth,height=2.5cm, keepaspectratio=false]{ENSG00000184009proportions_plot_no_title.pdf}        
     %  \par\vspace{-1em}%
       
     %  \vspace{-5mm}  % Negative space to pull images closer
    
        {\llap{\raisebox{\dimexpr\height-2mm}{\textbf{\tiny b}}}
         \vspace{-1.6mm}
    \tiny Transcript counts (ACTG1)
          \includegraphics[width=\linewidth,height=2.5cm, keepaspectratio=false]{ENSG00000184009cpm_plot_no_title.pdf}   
}     
          \end{subfigure}    
    \begin{subfigure}[t][6.05cm][t]{0.5\textwidth}
        \centering
           \llap{\raisebox{0mm}{\textbf{\tiny c }}}%
               \vspace{-2mm}
         \tiny Novel transcript in PCR-cDNA not detected in direct RNA samples  (ACTG1)
        \includegraphics[width=\textwidth,height=5.3cm, keepaspectratio=false]{ENSG00000184009coverage_plot_A549novel_exitron_log_scale.png}    
    \end{subfigure}
    
          \vspace{-5mm}
          
      \begin{subfigure}{0.25\textwidth}
             \centering  
                 \llap{\raisebox{0mm}{\textbf{\tiny d }}}%
               \vspace{-2mm}
           \tiny DTU genes with  a novel transcript
%\llap{\raisebox{\dimexpr\height-0.8em}{\textbf{\scriptsize a}} \scriptsize {Numbers of DTU}\hspace{0.01em}}%%
       \includegraphics[width=\linewidth,height=3.5cm, keepaspectratio=false]{/Users/rotem/Code/GitHub/sgnex-compare/code/analysis-pages/figures/figure4/dtu_genes_with_novel.png}
    \end{subfigure}
 \begin{subfigure}{0.65\textwidth}
        \centering
           \llap{\raisebox{0mm}{\textbf{\tiny e }}}%
               \vspace{-2mm}
    \tiny Structural differences detected between novel in PCRcDNA and dominant in direct RNA 
        \includegraphics[width=\linewidth,height=3.5cm, keepaspectratio=false]{events_numbers_in_novel_pairs.png}
    \end{subfigure}

%     
%    \begin{subfigure}{0.33\textwidth}
%        \centering
%        \llap{\raisebox{\dimexpr\height-0.8em}{\textbf{\scriptsize h}}}%
%        \includegraphics[width=\linewidth,height=3.5cm, keepaspectratio=false]{sqanti3_label_numbers_in_switching_pairs.png}
%    \end{subfigure}
%         
    

\end{figure}
\setcounter{figure}{3} 
  \captionof{figure}{\textbf{‘Novel’ transcript only observed in PCR cDNA are structurally different to the most dominant transcript in direct RNA.} \\
  a-c) Example of a DTU gene (ACTG1), where a ‘novel’ transcript was detected in PCR cDNA.  a) Estimated proportions in direct RNA (left) and PCR cDNA (right), estimated using DRIMSeq from ‘bambu’ raw transcript counts for the same gene. b) Transcript counts (log2cpm) obtained by ‘bambu’ for this gene. c) When compared to the dominant transcript in direct RNA (black), the novel transcript detected in PCR cDNA (dark blue) has a 5’ truncated  exon (can be seen on the right), and an inserted intron on the left. In this case, coverage plots corroborate this junction in all RT protocols, but not in direct RNA. 
d) Fraction of DTU genes detected, with the fraction of genes with novel isoforms in black.  e) Number of structural differences detected between transcript novel in cDNA and the one dominant in RNA. }
\end{center}
\end{document}
